\documentclass[11pt,a4paper]{article}
\usepackage[margin=25mm]{geometry}
\usepackage{amsmath,amssymb,amsthm,hyperref}
\setlength{\emergencystretch}{2em}
\newtheorem{lemma}{Lemma}
\newtheorem{theorem}[lemma]{Theorem}
\newcommand{\gen}[1]{\langle#1\rangle}
\title{Closing the paired interior between two singleton ends}
\author{Human Compression project: L044 continuation}
\date{}
\begin{document}\maketitle
Use all hypotheses and notation of the opposite-singleton budget
packet: $G\le A_n$, $n=2r\ge14$, has a connected transitive string
C-marking, all maximal omissions are intransitive, and
\[
 G_0=\operatorname{Alt}(\Omega\setminus\{a\}),\qquad
 G_{r-1}=\operatorname{Alt}(\Omega\setminus\{d\}).
\]
Every colour $2,\ldots,r-2$ has two ambient-eligible edges.
We additionally assume the alternating rank bound at strictly smaller
degrees at least twelve. The complete accepted primitive, anchored
imprimitive and small-tail bounds remain paid.
Write $\sigma=\rho_1$, $\eta=\rho_{r-1}$,
$H_C=G_{\{2,\ldots,r-2\}}$, $K_C=G_{\{3,\ldots,r-2\}}$,
and $K=G_{>2}$, as in the predecessor.

\begin{lemma}[The actual two short heads]
After renaming points, one of the following two literal pairs occurs,
with both generators fixing every point not displayed:
\[
\begin{array}{c|cc}
 &\rho_0&\rho_1\\ \hline
\text{five-point head}&(a\ b)(u\ v)&(b\ u)(v\ w)\\
\text{six-point head}&(a\ b)(u\ v)&(b\ x)(u\ w).
\end{array}                                             \tag{1}
\]
The first pair generates the natural five-point $D_{10}$; the second
generates $S_3$ diagonally on two three-point orbits.
The group $K$ fixes the displayed five or six points, respectively.
Moreover colour $1$ has at least two ambient-eligible edges.
\end{lemma}
\begin{proof}
We use the proved cases of the opposite-singleton budget, not an
unproved stronger interpretation of its statement. That proof
excludes $f_C=3$ and allows only $f_C=4$ or its literal residual
$f_C=2$. Here $f_C$ counts the $H_C$-fixed points in the
$G_{r-1}$-domain.

If $f_C=4$, then $H_C$ fixes $a,b,u,v$ and is transitive on
$V=\Omega\setminus\{a,b,u,v,d\}$. Since
$C_0=\gen{\sigma,H_C}$ is transitive on its complement of $a,d$,
the involution $\sigma$ must send each of $b,u,v$ into $V$.
A fixed singleton or a transposition among these three would instead
be a proper $C_0$-orbit. Their three images are distinct.
The subgroup $K_C$ fixes $a,b,u,v$ and their images under $\sigma$,
since it centralizes $\sigma$. Together with $d$ this gives eight
ambient fixed points. Its $r-4$ colours all have two eligible edges,
confined to at most $n-8=2(r-4)$ points. The strict localized interval
bound contradicts this, retaining fixed singleton orbits as usual.
Thus $f_C=4$ is impossible.

Consequently the predecessor's $f_C=2$ residual holds:
$E=\{u,v\}$ is a transitive $H_C$-orbit, $V$ is its other nontrivial
orbit of odd size $n-5$, and $\rho_0=(a\ b)(u\ v)$.
Its proof further establishes that $K_C$ fixes $E$ pointwise,
that $\eta$ fixes $E$ pointwise, and that
$x=b^\sigma\ne b$ lies in $E\cup V$.
If $x\in E$, name it $u$. Then $u^\sigma=b$ and
$v^\sigma=w\in V$, exactly as proved there. This gives the first
pair in (1), initially with a possible extra involution $\epsilon$
in $\sigma$ outside the displayed five points.

If $x\in V$, then $E^\sigma\ne E$ by transitivity of $C_0$.
It cannot be disjoint from $E$: in that case $K_C$ fixes the seven
local points $a,b,x,E,E^\sigma$, and the same strict bound on its
$r-4$ colours is impossible. Their intersection is therefore a
singleton. Since $\sigma$ is an involution, that singleton is fixed
by $\sigma$; name it $v$. Thus $u^\sigma=w\in V\setminus\{x\}$.
This is the second pair in (1), again initially decorated by an
exterior involution $\epsilon$.

In either case the two displayed transpositions of $\sigma$ are
even, so $\epsilon$ is even. In the first case
$(\rho_0\sigma)^5=\epsilon$, and in the second
$(\rho_0\sigma)^3=\epsilon$. The displayed product is respectively
a five-cycle or two disjoint three-cycles. Hence
\[
 \epsilon\in\gen{\rho_0,\rho_1}\cap G_0=\gen{\rho_1},
\]
where membership in $G_0$ follows from its full natural alternating
action and $\epsilon$ fixing $a$. The decoration fixes $b$ whereas
$\rho_1$ moves it, so $\epsilon=1$. This proves the exact marked
heads without quotient IP or a finite classification.
The predecessor's fixed-point argument shows that $K$ fixes all
five or six displayed points.

Finally $H_C$ preserves $E$ and $V$, while $\eta$ fixes $E,a,b$
and moves $d$ into $V$. Thus the actual omission
$G_1=\gen{\rho_0,H_C,\eta}$ preserves the three nonempty disjoint
sets $\{a,b\}$, $E$, and $V\cup\{d\}$.
It has at least three orbits. A single eligible $1$-edge could
join at most two of them when $\rho_1$ is adjoined, contradicting
transitivity of $G$. Colour $1$ is therefore paired too.
\end{proof}

\begin{lemma}[Both trimmed tails have an orbit deficit]
The group $K=G_{>2}$ is intransitive on its nonempty support and
fixes at least five points. The same holds for $G_{<r-3}$.
\end{lemma}
\begin{proof}
Suppose $K$ is transitive on its support $X$, of size $m$.
It has rank $e=r-3\ge4$, an actual even string C-marking and a
connected diagram. Its local maximal omissions are intransitive:
if one were transitive on $X$, adjoining the omitted generator,
which is supported on $X$, could not change the global omission
orbits, contradicting global transitivity.
All generators except $\rho_2$ preserve $X$, since the earlier
ones centralize $K$. Thus $\rho_2$ takes some point of $X$ outside,
and that point is fixed by all marks of $K$ except its first.
This is the exact attachment anchor required by the accepted bounds.

The complete accepted small-tail lemma leaves only the five-point
$D_{10}$ or six-point $A_5$ markings when $m<12$ and
$e\ge(m-1)/2$. Their ranks are at most three, excluded here.
For $m\ge12$, the accepted proper-primitive and anchored-imprimitive
bounds, and the explicit smaller-degree induction when
$K=\operatorname{Alt}(X)$, give
\[
                              2e\le m-1.
\]
The same inequality holds in the small case just discussed.
Since Lemma 1 supplies at least five fixed points,
$m\le n-5=2e+1$. Equality is therefore forced throughout:
$m=n-5$, and the fixed set consists of exactly five points.
The six-point head is excluded. In the five-point case the two
actual $G_2$-orbits are its head $D$ and $X$, because
$G_2=\gen{\rho_0,\rho_1,K}$ has these disjoint transitive factors.
The generator $\rho_0$ fixes $X$ and has just one fixed point,
$w$, in $D$. As $\rho_2$ commutes with $\rho_0$, any $2$-edge
between $D$ and $X$ must have endpoint $w$. There is at most one
such edge, contradicting the assumed two eligible $2$-edges.
Thus $K$ is intransitive on its support.

Lemma 1 proved that every interior colour $1,\ldots,r-2$ is
paired. Apply the whole argument in reversed order; all its precise
hypotheses hold. It gives the assertion for $G_{<r-3}$.
\end{proof}

\begin{theorem}[Two singleton ends with a paired interior]
There is no group satisfying the stated hypotheses.
\end{theorem}
\begin{proof}
Apply the accepted orbit-incidence deficit theorem with
$h=2$ and $g=r-3$. Since $r\ge7$, we have $g>h$.
The two groups $R=G_{>h}$ and $P=G_{<g}$ are intransitive on
their supports by Lemma 2, and each fixes at least five points.
Thus both required outside budgets hold:
\[
 2h=4\le n-|\operatorname{supp}(R)|-1,\qquad
 2(r-1-g)=4\le n-|\operatorname{supp}(P)|-1.
\]
Every middle colour $h<j<g$ is paired. With $p,q\ge2$ the
numbers of their nontrivial orbits, the complete incidence theorem
gives $2r\le n+3-p-q\le n-1$, contradicting $n=2r$.
\end{proof}

\paragraph{Attribution and exact remaining source boundary.}
This closes the stated subcase of Cameron--Fernandes--Leemans--Mixer,
\emph{Highest rank of a polytope for $A_n$}, arXiv:1605.09173v1,
Proposition 5.32, physical pages 44--45. The source's two-end
fixed-support and overlap strategy is retained. All accepted odd
singleton, opposite-end lift, selected-localization, anchored-tail
and orbit-incidence proofs remain paid. Short-head extraction uses
actual IP and the natural alternating omission. Other heads,
unpaired intermediate crossings, equality diagrams and remaining
rigidity/table obligations are separate. No new census, formalization,
original or question credit follows.
\end{document}
